% =============================================================
%  Section 3 -- Conceptual Framework
%  Label: sec:framework
% =============================================================

\section{Conceptual Framework}\label{sec:framework}

This section derives a rational benchmark for the household
forecast error and shows how diagnostic expectations generate
overshooting.
The model is deliberately parsimonious: its purpose is to
discipline the empirical test, not to estimate structural
parameters.
Full derivations appear in Appendix~\ref{app:model}.

\subsection{Setup}\label{sec:model_setup}

A household in province $p$ at survey wave $t$ observes the
cumulative meat-price change $s_{pt}$ between waves $t-1$ and $t$.
The household must forecast headline CPI inflation
$\pi_{p,\,t \to t+1}$ over the subsequent inter-wave period.

Meat enters headline CPI with weight $\omega \approx 0.05$.
The meat-price shock has persistence $\rho \in [0,1]$: a fraction
$\rho$ of the current shock carries forward into the next period's
headline inflation.
Formally, the headline CPI response to a meat shock of size $s$ is
\begin{equation}\label{eq:passthrough_model}
  \E[\pi_{p,\,t \to t+1} \mid s_{pt}] = \bar{\pi} + \omega\,\rho\,s_{pt},
\end{equation}
where $\bar{\pi}$ is unconditional expected inflation.
The household does not know $\rho$ with certainty but holds a
prior belief $\hat{\rho}$.

I focus on persistence uncertainty rather than weight uncertainty
because CPI weights are approximately stable across the sample
period, so persistent disagreement about $\omega$ would produce
a constant expectation bias absorbed by fixed effects.
Persistence uncertainty, by contrast, interacts with the size
of the shock and generates a coefficient on $s_{pt}$ in the
forecast-error regression.

\subsection{Rational Benchmark}\label{sec:rational_benchmark}

Under rational expectations, the household forecasts
$\E_t^{R}[\pi_{p,\,t \to t+1}] = \bar{\pi}
  + \omega\,\hat{\rho}\,s_{pt}$,
where $\hat{\rho}$ is the household's best estimate of
persistence.
The forecast error is
\begin{equation}\label{eq:rational_fe}
  \mathrm{FE}^{R}_{pt}
  = \pi_{p,\,t \to t+1} - \E_t^{R}[\pi_{p,\,t \to t+1}]
  = \omega(\rho - \hat{\rho})\,s_{pt} + \eta_{pt},
\end{equation}
where $\eta_{pt}$ collects shocks to headline CPI that are
orthogonal to the meat signal.
If the household overestimates persistence ($\hat{\rho} > \rho$),
it expects more meat-driven inflation than actually occurs, and
the forecast error is negative in provinces with positive shocks.
If the household underestimates persistence, the forecast error
is positive in high-shock provinces.

The coefficient on $s_{pt}$ in a regression of the forecast
error on the meat shock is therefore bounded:
\begin{equation}\label{eq:rational_bound}
  \beta_3^{R} = \omega(\rho - \hat{\rho}).
\end{equation}
Even in the worst case---where households believe the shock is
permanent ($\hat{\rho} = 1$) but it is fully transitory
($\rho = 0$)---the rational forecast-error coefficient is
$\beta_3^{R} = -\omega \approx -0.05$.
Under any intermediate belief, $|\beta_3^{R}| < \omega$.
The rational benchmark therefore predicts
$|\beta_3| \leq 0.05$ per unit of the meat shock.

This bound is the central object of the empirical test.
The number $0.05$ is not a free parameter---it is the CPI
weight on meat, determined by the NBS basket construction.
Any rational household, regardless of its belief about
persistence, generates a forecast-error coefficient no larger
than the mechanical weight of meat in the index, because
persistence governs the \textit{duration} of the meat
contribution, not its \textit{size per period}.

\subsection{Diagnostic Extension}\label{sec:diagnostic_extension}

Under diagnostic expectations \citep{Bordalo2018}, the household
overweights the state that recent signals make more
representative.
A large positive meat-price shock makes high-inflation states
more salient, so the household inflates its forecast by a factor
$\theta > 0$ beyond the rational update:
\begin{equation}\label{eq:diagnostic_forecast}
  \E_t^{\theta}[\pi_{p,\,t \to t+1}]
  = \bar{\pi}
    + \omega\,\hat{\rho}\,s_{pt}
    + \theta\,\omega\,s_{pt}.
\end{equation}
The additional term $\theta\,\omega\,s_{pt}$ arises because
the household tilts its subjective probability distribution
toward states with high likelihood ratios.
When $\theta = 0$, the household reverts to the rational
benchmark; when $\theta > 0$, it places too much probability
mass on high-inflation outcomes after observing a positive
meat-price shock.

The diagnostic forecast error is
\begin{equation}\label{eq:diagnostic_fe}
  \mathrm{FE}^{\theta}_{pt}
  = \omega(\rho - \hat{\rho})\,s_{pt}
    - \theta\,\omega\,s_{pt}
    + \eta_{pt},
\end{equation}
so the regression coefficient becomes
$\beta_3^{\theta} = \omega(\rho - \hat{\rho} - \theta)$.
With $\theta > 0$, the forecast error is more negative than under
rationality for any belief $\hat{\rho}$ and any true persistence
$\rho$.

The forecast-error regression is a sharper test than the
expectation regression.
The expectation coefficient $\beta_1$ depends on
$\omega(\hat{\rho} + \theta)$, so distinguishing the rational
component from the diagnostic component requires an independent
estimate of $\hat{\rho}$, which is not directly observed.
The forecast-error coefficient $\beta_3$, by contrast, is
bounded by $\omega$ under rationality regardless of
$\hat{\rho}$.
Any estimate of $|\hat{\beta}_3|$ exceeding $0.05$ rejects
rationality without requiring knowledge of the household's
belief about persistence.

\subsection{Testable Implications}\label{sec:testable}

The framework generates three predictions that the data can
discriminate.

\textit{Prediction 1: Bounded forecast errors under rationality.}
Under rationality, the forecast-error coefficient satisfies
$|\beta_3| \leq \omega \approx 0.05$.
If the estimated $\hat{\beta}_3$ exceeds this bound in
absolute value, rationality is rejected regardless of the
household's belief about persistence.
The bound is conservative: it assumes the worst-case belief
($\hat{\rho} = 1$ when $\rho = 0$); under any more accurate
belief, the rational coefficient is smaller.
In the data, the estimated coefficient is $\hat{\beta}_3 =
-6.07$, which exceeds the rational bound by a factor of 120.

\textit{Prediction 2: Excess expectation response.}
Under diagnostic expectations the expectation response
exceeds the rational update:
$\hat{\beta}_1 > \omega\,\hat{\rho}$.
The expectation result complements the forecast-error test
by showing \textit{where} the excess forecast error comes
from: through the numerator (expectations are too high) rather
than the denominator (realised CPI is too low).
In the preferred specification, $\hat{\beta}_1 = 0.242$
(s.e. = 0.174), which is positive but imprecisely
estimated due to the ordinal nature of the expectation variable.

\textit{Prediction 3: Commodity asymmetry.}
A household that overweights salient signals should respond more
to meat shocks (high purchase frequency, high media coverage)
than to grain shocks of comparable CPI weight.
The diagnostic model predicts that $\theta$ is larger for
meat than for grain because the purchase-frequency channel
amplifies the representativeness of meat-price signals.
Grain prices in China are stabilised by government procurement
and strategic reserves, so they vary less and attract less
household attention.
The placebo test in Section~\ref{sec:mechanism} confirms this
asymmetry: grain shocks produce near-zero forecast errors.

\subsection{Mapping to Empirical Equations}\label{sec:model_mapping}

The three predictions map directly to the three-equation system
estimated in Section~\ref{sec:strategy}.
Equation~1 regresses household expectations on the meat shock
and tests whether $\hat{\beta}_1 > 0$.
Equation~2 regresses forward realised headline CPI on the
meat shock and estimates the actual pass-through
$\omega\,\rho$.
Equation~3 regresses the forecast error
(realised CPI minus expectation) on the meat shock and tests
whether $|\hat{\beta}_3| > \omega$.

The empirical strategy uses province-by-wave fixed effects in
the preferred specification.
In the model, these fixed effects absorb $\bar{\pi}$ and any
province-by-wave shocks common to all households in a given cell.
The identifying variation is the deviation of a province's
meat shock from the region-by-wave mean.
The model predicts that the forecast-error bound
$|\beta_3| \leq 0.05$ holds even in this demanding
specification because the bound derives from the CPI weight,
not from the level of the shock.

Section~\ref{sec:calibration} compares the estimated
$\hat{\beta}_3$ against the rational bound calibrated
to the actual persistence of meat-price shocks in the data.

